\documentclass{beamer}

\usetheme{metropolis}

\usepackage[ngerman]{babel}
\usepackage[autostyle=true,german=quotes]{csquotes}
\usepackage[linewidth=1pt]{mdframed}
\usepackage{hyperref}
\usepackage{makecell}
\usepackage{pifont}
\usepackage{tikz}
\usetikzlibrary{positioning, calc, arrows, fit, decorations.pathreplacing, shapes, shapes.multipart, snakes, trees, automata}
\usepackage{verbatim}
\usepackage{textcomp}
\usepackage{centernot}
\usepackage{tabularx}
\usepackage[normalem]{ulem}
\usepackage{mathpartir}
\usepackage{pdfpages}
\usepackage{wasysym}
\usepackage{qtree}
\usepackage{pgfplots}

\batchmode

\hypersetup{
	colorlinks,
	urlcolor=blue,
	linkcolor=black % for ToC
}
\newenvironment{qaa}[1]{
	#1

	\begin{mdframed}
		\small
}{
	\end{mdframed}
}

\newcommand{\true}{\ding{51}}
\newcommand{\false}{\ding{55}}
\newcommand{\code}[1]{
	\begin{mdframed}
		\verbatiminput{#1}
	\end{mdframed}
}

% lambda calculus stuff

\newcommand{\app}[2]{#1\;#2}
\newcommand{\appr}[2]{\app{\uline{#1}}{\uwave{#2}}}
\newcommand{\lam}[2]{\lambda #1.\, #2}
\newcommand{\appp}[3]{\app{\app{#1}{#2}}{#3}}
\newcommand{\apppp}[4]{\app{\app{\app{#1}{#2}}{#3}}{#4}}
\newcommand{\apppr}[3]{\app{\appr{#1}{#2}}{#3}}

\newcommand{\lamlet}[3]{\texttt{let} \;\; #1 = #2 \;\; \texttt{in} \;\; #3}

\newcommand{\ctrue}{c_\text{true}}
\newcommand{\cfalse}{c_\text{false}}
\newcommand{\cand}{\text{AND}}

\newcommand{\csucc}{\text{succ}}
\newcommand{\cn}[1]{c_{#1}}

\newcommand{\cpair}{\text{pair}}
\newcommand{\cfst}{\text{fst}}
\newcommand{\csnd}{\text{snd}}
\newcommand{\ccurry}{\text{curry}}
\newcommand{\cuncurry}{\text{uncurry}}

\newcommand{\cnil}{\text{nil}}
\newcommand{\ccons}{\text{cons}}
\newcommand{\chead}{\text{head}}
\newcommand{\ctail}{\text{tail}}
\newcommand{\creplicate}{\text{replicate}}

\newcommand{\subst}[3]{(#1)\left[#2\,\to\,#3\right]}

% unification etc.

\newcommand{\functor}[2]{\texttt{#1}(#2)}
\newcommand{\pa}[1]{\texttt{#1}}
\newcommand{\plsunify}{\stackrel{!}{=}}
\newcommand{\unifier}[1]{\left[#1\right]}
\newcommand{\su}[2]{#1 \; \text{\pointer} \; #2}


\title{Tutorium 09: Prolog}
% \subtitle{}
\author{Paul Brinkmeier}
\institute{Tutorium Programmierparadigmen am KIT}
\date{10. Januar 2023}


\begin{document}

\begin{frame}
    \titlepage
\end{frame}

\section{Rückblick}

\begin{frame}{Rückblick}
    Bisherige Themen:

    \begin{itemize}
        \item Haskell: Funktionale Programmierung
        \item Lambda-Kalkül: Beta-Reduktion, Church-Encodings
        \item Typisierung: Lambda-Typen, Let, Typinferenz
        \item Prolog: Logische Programmierung
    \end{itemize}

    Ab nächster Woche:

    \begin{itemize}
        \item Parallelprogrammierung: MPI, Java
        \item Design by Contract: OpenJML
        \item Compiler: Parser + Java ByteCode
    \end{itemize}
\end{frame}

\begin{frame}{Klausuranmeldung}
    \begin{itemize}
        \item Klausur: 31.03.2023 um 11:30
        \item Anmeldung ab Montag bis 24.02.2023 im \href{https://campus.studium.kit.edu/exams/registration.php?pguid=0x87BA4E33FD4947E8ADAB43370BDFA852}{CAS}
        \item Dauer: 2 Stunden
        \item Erlaubtes Hilfsmaterial: \enquote{Alles aus Papier}
    \end{itemize}
\end{frame}

\section{Übungsblätter}

\section{Prolog}

\subsection{Kleine Aufgaben}

\begin{frame}{Maximum-Funktion}
    Schreibt ein Prädikat \texttt{max/3}, das das Maximum bestimmt:

    \code{code/max.out}

    \begin{itemize}
        \item Schreibt die Funktion zuerst in Haskell
        \begin{itemize}
            \item Verwendet die Guard-Notation
            \item Wie viele Vergleiche benötigt man?
        \end{itemize}
        \item Lässt sich die Haskell-Version 1:1 übersetzen?
    \end{itemize}
\end{frame}

\begin{frame}{Einteilen in Gruppen}
    Schreibt ein Prädikat \texttt{classify/2}, das Zahlen in folgende Gruppen einteilt:

    \begin{itemize}
        \item \texttt{negativ}: Zahlen kleiner als 0
        \item \texttt{interessant}: Zahlen von 0 bis 42 (inklusiv)
        \item \texttt{zugross}: Zahlen über 42
    \end{itemize}

    \code{code/classify.out}
\end{frame}

\subsection{Cut}

\begin{frame}{Cut!}
    \only<1>{
        \code{../demos/classify1.pro}

        \begin{itemize}
            \item Wenn die Zahl \texttt{negativ} ist, können \texttt{X >= 0} und \texttt{X > 42} nicht stimmen
            \item $\leadsto$ Prolog sollte diese eigentlich überspringen
            \item Prolog macht aber keine Arithmetik für uns
            \item Aushilfe: Der Cut\texttt{!}
        \end{itemize}
    }

    \only<2>{
        \code{../demos/classify2.pro}

        \begin{itemize}
            \item Vermeidung redundanter Berechnungen: Cut, geschrieben \texttt{!}
            \item Cut kann immer einmal erfüllt werden
            \item Reerfüllung lässt aber die ganze Regel fehlschlagen
            \item Verschiedene Erklärungen:
            \begin{itemize}
                \item Cut \enquote{schneidet Reerfüllungsbaum ab}
                \item Cut entfernt Choice Points
                \item Cut \enquote{verpflichtet} Prolog zur Erfüllung einer bestimmten Regel
            \end{itemize}
        \end{itemize}
    }

    \only<3>{
        \code{../demos/classify2.pro}

        \begin{itemize}
            \item Schließen sich die Tests einzelner Fälle gegenseitig aus, kann man nach ihnen gefahrlos Cuts einfügen
            \item \enquote{Grüner} Cut: Ändert nicht Verhalten, nur Laufzeit des Programms
            \item Hier:
            \begin{itemize}
                \item \texttt{X < 0} $\leadsto$ \texttt{X >= 0} und \texttt{X > 42} unmöglich
                \item \texttt{X =< 42} $\leadsto$ \texttt{X > 42} unmöglich
                \item \texttt{!} $\Leftrightarrow$ \enquote{Ab hier kann kein anderer Fall mehr eintreten, fertig}
            \end{itemize}
        \end{itemize}
    }

    \only<4>{
        \code{../demos/classify3.pro}

        \begin{itemize}
            \item Einige Tests lassen sich \emph{durch Cuts ersetzen}
            \item \enquote{Roter} Cut: Ändert Verhalten \emph{und} Laufzeit des Programms
            \item Fehleranfällig! Wenn möglich vermeiden.
            \item Programme ohne rote Cuts sind
            \begin{itemize}
                \item einfacher zu verstehen.
                \item etwa genauso performant.
            \end{itemize}
            \item Faustregel: Reihenfolge der Regeln sollte austauschbar sein
        \end{itemize}
    }
\end{frame}

\begin{frame}{Summe Vorwärts und Rückwärts}
    Schreibt ein Prädikat \texttt{sum/3}, sodass \texttt{sum(A, B, C)} die Gleichung

    \begin{equation*}
        \texttt{A} + \texttt{B} = \texttt{C}
    \end{equation*}

    löst, wenn \emph{zwei oder mehr Variablen belegt sind}.

    \begin{itemize}
        \item 3 Fälle.
        \item Verwendet \texttt{number}, um zu Prüfen ob eine Variable mit einer Zahl belegt ist.
    \end{itemize}
\end{frame}

\subsection{Klausuraufgabe}

\begin{frame}{Klausur SS15: Haus vom Nikolaus}
\end{frame}

\section{Boggle}

{
\usebackgroundtemplate{\includegraphics[width=\paperwidth]{images/boggle.jpg}}
\begin{frame}[plain]
\end{frame}
}

\begin{frame}{Boggle}
    \begin{figure}
        \includegraphics[width=0.6\textwidth]{images/boggle.jpg}
    \end{figure}
    \only<1>{
    \begin{itemize}
        \item 16 Buchstabenwürfel werden gleichzeitig gewürfelt
        \item Jedes gefundene Wort gibt Punkte
        \begin{itemize}
            \item Wort: \emph{Mindestens drei zusammenhängende Würfel}
            \item Regeln ähnlich wie Scrabble: Wörterbuch, keine Abkürzungen/Eigennamen, etc.
        \end{itemize}
    \end{itemize}
    }
    \only<2>{
        \begin{itemize}
            \item Schreibt ein Programm in Prolog, das ein Boggle löst
            \item Vorlage: \texttt{demos/boggle.pro}
            \begin{itemize}
                \item Implementiert: \texttt{matchCharAt}, \texttt{neighbor} und \texttt{search}
                \item Das vorgegeben Prädikat \texttt{word} definiert unsere Wortliste
            \end{itemize}
        \end{itemize}
    }
\end{frame}

\end{document}
